% [原文 p.115]
\begin{proof}[证明（续）]
取 $(p,q)\in I$，使 $aq-rp$ 达到最大值，并将其值记为 $d/v$。选取步骤 7 中那样的奇点消解 $f$。在带状区域 $0<aq'-rp'<d/v$ 上，与以前同样可得下式：
\[
\chi\bigl(F,\left.(f^*(p'H+q'a(K_X+\Delta))+\lceil A'\rceil)\right|_F\bigr)
=h^0\bigl(F,\left.(f^*(p'H+q'a(K_X+\Delta))+\lceil A'\rceil)\right|_F\bigr).
\]
由步骤 3 可知，在取 $\epsilon=1$ 的带状区域 $0<aq'-rp'<1$ 中，存在满足 $h^0(F,\ldots)\neq 0$ 的点 $(p',q')$。然而此时 $aq'-rp'<d/v=aq-rp$ 自动成立。矛盾与步骤 7--9 同样导出。至此，有理性定理的证明结束。
\end{proof}

\section{非消没定理}
\takeaway{非消没定理断言：在 klt 与丰富性的假设下，$H^0(X,mD+\lceil G\rceil)$ 对 $m\gg 0$ 不为零。证明先化归到 $X$ 非奇异的情形，再以曲线的重数估计造出丰富的辅助除子，最后用消没定理与对维数的归纳完成。}

\begin{proof}
本节证明取自 [Sho85]。

\par\medskip\noindent\textbf{步骤 0}\quad 设 $f:X'\to X$ 是 $X$ 的射影奇点消解，并设 $G'$ 是满足 $K_{X'}-G'\equiv f^*(K_X-G)$ 的除子。此时 $(X',-G')$ 是 klt 的，且 $af^*D+G'-K_{X'}\equiv f^*(aD+G-K_X)$。与命题 2.61 同样地，选取小的有效除子 $F$，使 $af^*D+G'-K_{X'}-F$ 丰富，且 $(X',-G'+F)$ 是 klt 的。令 $G'':=G'-F$。此时 $f_*(G'')\le G$ 成立，故
\[
h^0(X',mf^*D+\lceil G''\rceil)\le h^0(X,mD+\lceil G\rceil).
\]
由此可知，定理 3.4 只需在下面的附加条件下证明即可。以后在证明中假定这一条件。
（$*$）$X$ 是非奇异射影簇，$(X,-G)$ 是 klt 的，且 $aD+G-K_X$ 对适当的 $a>0$ 是丰富的。

\par\medskip\noindent\textbf{步骤 1}\quad 由于 $(X,-G)$ 是 klt 的，$G$ 的各成分的系数都 $>-1$；换言之，$\lceil G\rceil$ 是有效的。若 $D$ 是数值平凡的，则
\[
h^0(X,mD+\lceil G\rceil)=\chi(X,mD+\lceil G\rceil)=\chi(X,\lceil G\rceil)=h^0(X,\lceil G\rceil)\neq 0.
\]
其中中间的等式由命题 2.57 得到，两端的等式由定理 2.47 得到。因此定理的结论显然成立。以后可以假定 $D$ 不是数值平凡的。
% [Tutor Note: 步骤 0 把一般情形化归为 $X$ 非奇异、$(X,-G)$ 为 klt 的情形；由于 $G''=G'-F\le G$，化归后的截面空间更小，故在那里证明非零性即已足够。]

% [原文 p.116]
\par\medskip\noindent\textbf{步骤 2}\quad 选取不在 $G$ 的支集中的点 $x\in X$。（在下面构造 $f$ 时，首先爆破这个点。）我们要证明：可以适当选取 $q_0\in\Z_{>0}$，使得对每个 $q\ge q_0$，存在 $\Q$-除子 $M(q)\equiv(qD+G-K_X)$ 满足 $\mathrm{mult}_xM(q)>2\dim X$。

为此令 $d=\dim X$。由于 $D$ 是 nef 的，对任何丰富除子 $A$ 都有 $(D^e\cdot A^{d-e})\ge 0$。由此
\[
((qD+G-K_X)^d)=(((q-a)D+aD+G-K_X)^d)\ge d(q-a)(D\cdot(aD+G-K_X)^{d-1}).
\]
适当选取满足 $(D\cdot C)>0$ 的曲线 $C$ 与 $p\in\Z_{>0}$，可使 $((p(aD+G-K_X))^{d-1})$ 由 $C$ 与有效 $1$-循环之和代表。由此 $(D\cdot(aD+G-K_X)^{d-1})>0$。因此，当 $q\to\infty$ 时右端 $\to\infty$。于是由 Riemann-Roch 的弱定理与 Serre 消没定理，
\[
h^0(e(qD+G-K_X))\ge\frac{e^d}{d!}(qD+G-K_X)^d+(\text{$e$ 的低次项}).
\]
另一方面，$M(q,e)$ 在 $x$ 处的重数 $>2de$，相当于关于 $M(q,e)\in|e(qD+G-K_X)|$ 的、至多 $\frac{e^d}{d!}(2d)^d+(\text{$e$ 的低次项})$ 个条件。由于当 $q\to\infty$ 时 $(qD+G-K_X)^d\to\infty$，截面的个数多于条件的个数。这样就找到了满足 $\mathrm{mult}_xM(q,e)>2de$ 的 $M(q,e)\in|e(qD+G-K_X)|$。令 $M(q):=M(q,e)/e$，即证得上述断言。

\par\medskip\noindent\textbf{步骤 3}\quad 设 $f=f(q):Y\to X$ 是 $(X,G+M(q))$ 的对数奇点消解，且它支配 $x\in X$ 处的爆破 $B_xX$。按如下方式引入记号。
\begin{enumerate}[label=（\arabic*）,leftmargin=2.8em,itemsep=1pt,topsep=2pt,parsep=0pt]
\item $K_Y+f^*G\equiv f^*K_X+\sum b_jF_j$，其中由假定 $b_j>-1$；
\item $\frac12f^*(aD+G-K_X)-\sum p_jF_j$ 对适当的 $0<p_j\ll 1$ 是丰富的；
\item $f^*M(q)=\sum r_jF_j$，其中 $F_0$ 是爆破 $x$ 所得的例外除子。
\end{enumerate}

% [原文 p.117]
\par\medskip\noindent\textbf{步骤 4}\quad 定义
\[
N(b,c):=bf^*D+\sum(-cr_j+b_j-p_j)F_j-K_Y.
\]
与以前同样，我们希望 $N(b,c)$ 是丰富的。为此作如下计算。
\[
\begin{aligned}
N(b,c)&=bf^*D+\sum(-cr_j+b_j-p_j)F_j-K_Y\\
&\equiv bf^*D-cf^*(qD+G-K_X)-\sum p_jF_j+f^*G-f^*K_X\\
&=(b-a-c(q-a))f^*D+(1-c)f^*(aD+G-K_X)-\sum p_jF_j\\
&=(b-a-c(q-a))f^*D\quad(\text{若 }b-a-c(q-a)\ge 0\text{，则 nef})\\
&\quad+\Bigl(\frac12-c\Bigr)f^*(aD+G-K_X)\quad(\text{若 }1/2\ge c\text{，则 nef})\\
&\quad+\frac12f^*(aD+G-K_X)-\sum p_jF_j\quad(\text{丰富}).
\end{aligned}
\]
由此可知，只要 $1/2\ge c$ 与 $b\ge a+c(q-a)$ 同时成立，$N(b,c)$ 就是丰富的。

\par\medskip\noindent\textbf{步骤 5}\quad 取 $c=\min\{(1+b_j-p_j)/r_j\}$，其中最小值取遍所有满足 $r_j>0$ 的 $j$。此时 $c>0$。与以前同样，先略微变动系数 $p_j$，使这个最小值只在唯一的 $j'$ 处达到，并令 $F=F_{j'}$。由 $x\notin\mathrm{Supp}\,G$ 可知 $b_0=d-1$、$r_0>2d$，故 $c<(1+(d-1)-p_0)/2d<1/2$。既然得到 $c<1/2$，则当 $b\ge a+c(q-a)$ 时 $N(b,c)$ 是丰富的。
% [Tutor Note: 步骤 5 中的 $c$ 由系数 $b_j,p_j,r_j$ 的最小值定义；$c<1/2$ 来自 $x\notin\mathrm{Supp}\,G$ 时的 $b_0=d-1$ 与 $r_0>2d$，这正是把 $N(b,c)$ 调成丰富除子的关键。]

\par\medskip\noindent\textbf{步骤 6}\quad 剩下部分与无基点定理、有理性定理的证明相同。由
\[
N(b,c)=bf^*D+A-F-K_Y
\]
定义 $\Q$-除子 $A$，并令 $f^*G=\sum g_jF_j$。若 $F_j$ 不是 $f$-例外的，则 $b_j=g_j$。$A$ 中 $F_j$ 的系数是 $-cr_j+b_j-p_j\ (<b_j)$，故
\[
\lceil A\rceil\le f^*\lceil G\rceil+(\text{$f$-例外除子}).
\]
在拉回上添加例外除子并不增加 $H^0$，因此
\[
H^0(Y,bf^*D+\lceil A\rceil)\subset H^0(Y,bf^*D+f^*\lceil G\rceil)=H^0(X,bD+\lceil G\rceil).
\]
由于 $N(b,c)$ 是丰富的，
\[
H^1(Y,bf^*D+\lceil A\rceil-F)=H^1(Y,bf^*D+\lceil A-F\rceil)=0.
\]
% [原文 p.118]
因此，$H^0(X,bD+\lceil G\rceil)\neq 0$ 只需证明
\[
H^0(F,\left.(bf^*D+\lceil A\rceil)\right|_F)\neq 0.
\]
最后的这个不等式由对 $\dim X$ 的归纳得到。事实上，若假定非消没定理对维数 $<\dim X$ 的簇已经证明，即可将它应用于 $F$，从而证得该不等式。
\end{proof}

\section{相对化}
\takeaway{相对化一节把无基点定理与锥定理搬到相对情形：对 $f:X\to Y$ 定义 $f$-nef、$f$-大、$f$-无基点，给出相对无基点定理 3.24 与相对锥定理 3.25，并在 $Y$ 为射影时用相对 Kleiman 判定法从绝对版推出。}

本节叙述无基点定理与锥定理的相对版本。它们往往比绝对版本更易于应用。

\par\medskip\noindent\textbf{3.21}\quad 下面是相对设定最原始的原型。
\begin{enumerate}[label=（\arabic*）,leftmargin=2.8em,itemsep=1pt,topsep=2pt,parsep=0pt,start=0]
\item $f:X\to Y$ 是特征 $0$ 的代数闭域上定义的射影簇之间的态射。此时，无基点定理与锥定理的相对版可借助相对 Kleiman 判定法（定理 1.44）从绝对版推出。本节最后将给出这一论证。
\end{enumerate}
更一般的相对版在两种情形下已知。
\begin{enumerate}[label=（\arabic*）,leftmargin=2.8em,itemsep=1pt,topsep=2pt,parsep=0pt,start=1]
\item $f:X\to Y$ 是特征 $0$ 的代数闭域上定义的固有态射。此时的证明只需对绝对情形的证明稍作改动，沿同一思路即可完成。详细处理请参阅 [KMM87]。
\item $f:X\to Y$ 是复解析流形之间的射影态射。此时必须假定 $Y$ 是紧子集的充分小邻域。这是为了避开 $X$ 是 $Y$ 的无穷多个点处的爆破这样的情形。这种情形在 [Nak87] 中处理。
\end{enumerate}
与定理 3.24 同样的结果，对 $f:X\to Y$ 是复解析流形的固有态射、且双有理于射影（或 Kähler）态射的情形也应当成立，但目前一无所知。这种情形下应当使用的消没定理在 [Tak94] 中得到证明。\hfill$\square$

\begin{definition}{3.22}
$f:X\to f(X)\subset Y$ 是满足条件 3.21(0)--(2) 之一的固有态射，$X$ 不可约，$n$ 是 $f$ 的一般纤维的维数。
% [原文 p.119]
$X$ 上的 Cartier 除子 $D$ 若对适当的 $c>0$，当 $k\gg 1$ 时成立
\[
\rank f_*\cO_X(kD)>c\cdot k^n,
\]
则称 $D$ 是 $f$-大的（$f$-big）（参见定义 2.59）。虽是显然之事，但须注意：若存在 $Y$ 的开集 $U$ 使 $f^{-1}(U)\neq\varnothing$ 且 $D|_{f^{-1}(U)}$ 是 $f|_{f^{-1}(U)}$-大的，则 $D$ 是 $f$-大的。

若 $f_*\cO_X(D)\neq 0$，则自然同态 $f^*f_*\cO_X(D)\to\cO_X(D)$ 定义 $Y$ 上的有理映射 $\phi_{|D|/Y}:X\dashrightarrow\PP_Y(f_*\cO_X(D))$。若自然同态 $f^*f_*\cO_X(D)\to\cO_X(D)$ 是满射，则称除子 $D$ 是 $f$-无基点的（$f$-free，原书作 $f$-自由）。此时 $\phi_{|D|/Y}$ 是态射。
\end{definition}
% [Tutor Note: $f$-大与 $f$-无基点分别是「大」「无基点」的相对化：前者用 $f_*\cO_X(kD)$ 增长阶的 $k^n$ 系数刻画，后者要求自然同态 $f^*f_*\cO_X(D)\to\cO_X(D)$ 为满射。]

下面叙述的 $f$-大性的性质与引理 2.60 同样证明。

\begin{lemma}{3.23}
$f:X\to Y$ 是射影的，$X$ 不可约，$D$ 是 $X$ 上的 Cartier 除子。此时下列条件等价。
\begin{enumerate}[label=（\arabic*）,leftmargin=2.8em,itemsep=1pt,topsep=2pt,parsep=0pt]
\item $D$ 是 $f$-大的。
\item 对任意拟射影的或 Stein 的开集 $U\subset Y$，存在 $m>0$ 使 $mD|_{f^{-1}(U)}\sim A+E$（$A$ 是 $f|_{f^{-1}(U)}$-丰富的，$E$ 是有效的）。
\item 存在 $m>0$ 使有理映射 $\phi_{|mD|/Y}$ 是双有理的。
\item 存在 $m>0$ 使 $\phi_{|mD|/Y}$ 的像与 $X$ 维数相等。
\item 对任意开集 $U\subset Y$ 与 $U$ 上的大除子 $M_U$，存在 $m>0$ 使 $D|_{f^{-1}(U)}+m(f|_{f^{-1}(U)})^*M_U$ 是大的。
\end{enumerate}
\end{lemma}

\begin{theorem*}[相对无基点定理（relative basepoint-free theorem）]{3.24}
$(X,\Delta)$ 是 klt 对，$\Delta$ 有效，且 $f:X\to Y$ 是满足条件 3.21(0)--(2) 之一的固有态射。$D$ 是 $f$-nef 的 Cartier 除子，且对适当的 $a>0$，$aD-K_X-\Delta$ 是 $f$-nef 且 $f$-大的。此时，对所有 $b\gg 0$，$bD$ 是 $f$-无基点的。
\end{theorem*}

\begin{theorem*}[相对锥定理（relative cone theorem）]{3.25}
$(X,\Delta)$ 是 klt 对，$\Delta$ 有效，且 $f:X\to Y$ 是满足条件 3.21(0)--(2) 之一的射影态射。此时，
\begin{enumerate}[label=（\arabic*）,leftmargin=2.8em,itemsep=1pt,topsep=2pt,parsep=0pt]
\item 存在至多可数个有理曲线 $C_j\subset X$ 满足下列条件。
\[
f(C_j)=(\text{点}),\qquad 0<-(K_X+\Delta)\cdot C_j\le 2\dim X,
\]
\[
\NE(X/Y)=\NE(X/Y)_{(K_X+\Delta)\ge 0}+\sum\R_{\ge 0}[C_j].
\]
% [原文 p.120]
\item 对任意 $\epsilon>0$ 与 $f$-丰富除子 $H$，
\[
\NE(X/Y)=\NE(X/Y)_{(K_X+\Delta+\epsilon H)\ge 0}+\sum_{\text{有限}}\R_{\ge 0}[C_j].
\]
\item 设 $F\subset\NE(X/Y)$ 是 $(K_X+\Delta)$-负的极锥面。此时存在 $Y$ 上射影的 $Z$ 与 $Y$-态射 $\mathrm{cont}_F:X\to Z$，使得 $(\mathrm{cont}_F)_*\cO_X=\cO_Z$，且对任意不可约曲线 $C\subset X$ 有
\[
\mathrm{cont}_F(C)=(\text{点})\iff[C]\in F.
\]
$\mathrm{cont}_F$ 的唯一性是显然的，它称为 $F$ 的收缩。
\item $F$ 与 $\mathrm{cont}_F:X\to Z$ 如 (3) 所述。$L$ 是 $X$ 上的线丛，且对任意满足 $[C]\in F$ 的曲线 $C$ 都有 $(L\cdot C)=0$。此时存在 $Z$ 上的线丛 $L_Z$ 使 $L\cong\mathrm{cont}_F^*L_Z$。
\end{enumerate}
\end{theorem*}

本节余下的部分，在 $f:X\to Y$ 是射影簇的态射的情形下证明相对版定理。
% [Tutor Note: 3.21 把相对版分成三种情形：$Y$ 为簇、$f$ 固有、$f$ 为复解析流形间的射影态射；后两种需要额外条件或另换文献，可见相对版并非绝对版的免费推广。]

\par\medskip\noindent\textbf{3.26}（情形 3.21(0) 中定理 3.24 的证明）\quad 把 $D$ 换成 $D+f^*(\text{丰富除子})$，并把 $\Delta$ 稍作变化，在保持 $(X,\Delta)$ 是 klt 这一条件的同时，使 $D$ 成为 nef 的，使 $aD-K_X-\Delta$ 成为丰富的。（一般地，$(f\text{-nef})+\nu f^*(\text{丰富})$ 未必对任何 $\nu>0$ 都是 nef 的（例 1.46）。）一旦做到这一点，就可以把绝对无基点定理 3.3 应用于 $D$，从而得到定理 3.24，这一点请注意。

由于 $aD-K_X-\Delta$ 是 $f$-大的，取 $Y$ 上适当的丰富除子 $H$，$a(D+f^*H)-K_X-\Delta$ 就是大的（引理 3.23）。存在 $X$ 上的有效 $\Q$-Cartier $\Q$-除子 $E$，使 $0<\epsilon\ll 1$ 时 $a(D+f^*H)-K_X-\Delta-\epsilon E$ 是丰富的（命题 2.61）。由推论 2.35(2)，$0<\epsilon\ll 1$ 时 $(X,\Delta+\epsilon E)$ 是 klt 的，故不妨把 $\Delta+\epsilon E$ 改名为 $\Delta$，把 $D+f^*H$ 改名为 $D$。这样，在保持其他条件的同时，可使 $F:=aD-K_X-\Delta$ 成为丰富的。

因此，由 $aD\equiv K_X+\Delta+F$ 与定理 3.7(2)，$\NE(X)$ 的极射线 $\R_{\ge 0}[\ell_j]$ 中包含于 $D_{<0}$ 的至多只有有限个。注意 $(\ell_j\cdot D)<0$ 且 $D$ 是 $f$-nef 的，故 $\ell_j$ 不含于 $f$ 的纤维中。
% [原文 p.121]
因此取 $m\gg 0$，就能使这些有限个 $\ell_j$ 都满足 $(\ell_j\cdot D+mf^*H)>0$。于是 $D+mf^*H$ 是 nef 的（定理 3.7(2)），$a(D+mf^*H)-K_X-\Delta$ 是丰富的，而这正是所要证明的。\hfill$\blacksquare$

\begin{propositionp}{3.27}
$f:X\to Y$ 是射影簇的态射。此时下列成立。
\begin{enumerate}[label=（\arabic*）,leftmargin=2.8em,itemsep=1pt,topsep=2pt,parsep=0pt]
\item $f$ 定义 $\R$-线性映射 $f_*:N_1(X)\to N_1(Y)$。
\item $f_*NE(X)\subset NE(Y)$，$f_*\NE(X)\subset\NE(Y)$。
\item $N_1(X/Y)\subseteq\ker f_*$，但等式一般并不成立。
\item $\NE(X/Y)=\NE(X)\cap N_1(X/Y)$。
\end{enumerate}
\end{propositionp}

\begin{proof}
(1) 已在定义 1.16 中确认；(2) 与 (3) 的前半是显然的。(2) 的后半由前半与 $f_*$ 的连续性导出。(3) 的后半由例 1.46 中 $\rho(X/Y)=\rho(Y)=1$、$\rho(X)\ge 3$ 得出。需要证明的只有 (4)。

首先「$\subset$」是显然的。假定 $\NE(X/Y)\subsetneq\NE(X)\cap N_1(X/Y)$ 并导出矛盾。存在 $N_1(X/Y)$ 上的线性函数 $H$，满足 $H_{>0}\supset\NE(X/Y)\setminus\{0\}$，且对适当的 $\xi\in\NE(X)\cap N_1(X/Y)$ 有 $(H\cdot\xi)<0$。略微变动 $H$ 的系数，可以假定 $H$ 是 $X$ 上的 $\Q$-Cartier 除子。

由相对 Kleiman 判定法（定理 1.44），$H$ 是 $f$-丰富的。由命题 1.45，适当选取 $Y$ 上的丰富除子 $M$，$H+f^*M$ 在 $X$ 上丰富。于是 $0<(H+f^*M\cdot\xi)=(H\cdot\xi)<0$，矛盾。至此 (4) 的证明完成。
\end{proof}

\par\medskip\noindent\textbf{3.28}（情形 3.21(0) 中定理 3.25 的证明）\quad 先证定理 3.25(1)。取 $\xi\in\NE(X/Y)$。由定理 3.7(1)，可用 $\eta\in\NE(X)_{(K_X+\Delta)\ge 0}$ 与 $r_j>0$ 表为 $\xi=\eta+\sum r_j[C_j]$。作用 $f_*$，得 $0=f_*\xi=f_*\eta+\sum r_jf_*[C_j]$。其中 $f_*\eta,f_*[C_j]\in\NE(Y)$（命题 3.27(2)）。由于 $Y$ 是射影的，$\NE(Y)$ 不含直线（定理 1.18）。因此 $f(C_j)$ 是点。也就是说，对每个 $j$ 有 $[C_j]\in\NE(X/Y)$，且 $\eta=\xi-\sum r_j[C_j]\in N_1(X/Y)$。于是 $\eta\in\NE(X)\cap N_1(X/Y)=\NE(X/Y)$。定理 3.25 的剩余部分很简单。\hfill$\blacksquare$
% [Tutor Note: 3.28 只用 (1) 的锥分解说明 $C_j$ 落在 $f$ 的纤维内；之所以能这样推出，是因为射影的 $Y$ 的曲线锥不含过原点的直线（定理 1.18）。]

% [原文 p.122]
\section{极小模型纲领的运行}
\takeaway{本节把极小模型纲领推广到对数对：先选取合适的对之类（$\Q$-分解的 dlt 对），再用锥定理与收缩定理驱动相对 MMP 的每一步，并指出其中哪些环节已经确立、哪些仍待解决。}

在 2.14 中也已概述了非奇异射影簇的极小模型纲领（MMP）。本节把它推广到对 $(X,\Delta)$，并说明本章证明的定理如何应用于极小模型纲领。

由于 MMP 或许不仅在 $\C$ 上的射影簇、而且在更多情形下都能运作，我们尽量在一般的设定下概述。

\par\medskip\noindent\textbf{3.29}（对之类别的选取）\quad $X$ 是域上的正规概形或正规复解析流形。为选取纲领中应当使用的对 $(X,\Delta)$ 的类别，请留意以下 3 个目标。
\begin{enumerate}[label=（\arabic*）,leftmargin=2.8em,itemsep=1pt,topsep=2pt,parsep=0pt]
\item 这一类别必须包含由非奇异 $X$ 与各系数不超过 $1$ 的有效 snc 除子 $\Delta$ 组成的对 $(X,\Delta)$。
\item MMP 的每一步都不应越出这一类别。
\item MMP 的每一步都应当存在。
\end{enumerate}
使这些成立的最大类别，多半是全体 lc 对的类别。(2) 与 (3) 即使在已知其为正确的情形，证明中也有困难之处。因此考虑更小的类别，即由 $X$ 为 $\Q$-分解的射影簇的全体 dlt 对 $(X,\Delta)$ 组成的类别。

这是使条件 (1) 与 (2) 成立的最小类别。（这由定理 2.44 较容易得到。）而 (3) 在维数 2、3 时已知。

$\Q$-分解性的假定，若从非奇异簇出发运行纲领，则非常自然，证明在许多情形下也变得简单。另一方面，它又是不稳定的条件。它关于 Euclid（或 étale）拓扑不是局部性质，在归纳证明中从簇转移到除子时，要确认它是否仍然成立非常困难。\hfill$\square$

\par\medskip\noindent\textbf{3.30}（对的双有理几何）\quad 一旦确定了应当使用的对之类，下一个必须直面的问题就是：

两个对 $(X_i,\Delta_i)$ 何时是双有理的？

当然，$X_1$ 与 $X_2$ 应当是双有理的。$\phi:X_1\dashrightarrow X_2$ 是双有理映射。
% [原文 p.123]
此外还必须确立 $\Delta_1$ 与 $\Delta_2$ 的关系。最容易的或许是令 $\Delta_2=\phi_*\Delta_1$。若 $\phi$ 是态射，或更一般地 $\phi^{-1}$ 不含例外除子，则这看起来是妥当的条件。然而，若 $E\subset X_2$ 是 $\phi^{-1}$ 的例外除子，则没有特别理由把 $E$ 的系数取为 $0$。事实上，从种种观点看，把 $\phi^{-1}$ 的一切例外除子的系数都取为 $1$，反而更为自然。

对于双有理等价的对，目前还不知道什么是好的定义。即使 $\phi:X_1\to X_2$ 是态射、$\Delta_2=\phi_*\Delta_1$ 的情形，也不应把两个对称为「双有理」的。这一点从下面的例子也可看出。

设 $X_2=\PP^2$，$\Delta_2=(x_0x_1x_2=0)$。设 $\phi:X_1\to X_2$ 是在 $(1:0:0)$ 处的爆破，并令 $\Delta_1=\phi_*^{-1}\Delta_2$。此时 $\Delta_2=\phi_*\Delta_1$。
% [Tutor Note: 此处 $\phi_*^{-1}\Delta_2$ 指 $\Delta_2$ 沿爆破 $\phi$ 的严格变换（proper transform），并非 $\phi_*$ 的逆映射；下文「$\phi^{-1}$ 的例外除子」也是按拉回像理解。]

计算 $K+\Delta$ 的倍数的全局截面，则对任何 $m>0$ 都有
\[
h^0(X_2,\cO_{X_2}(m(K_{X_2}+\Delta_2)))=1,\qquad
h^0(X_1,\cO_{X_1}(m(K_{X_1}+\Delta_1)))=0.
\]
如这个例子所示，目前作为对的极小模型的定义，尚不知道有值得一试的定义。在确立了对的极小模型纲领的框架之后，再定义适合下一节使用的极小模型概念，或许更好。\hfill$\square$
% [Tutor Note: 3.30 的 $\PP^2$ 例子表明：即使 $\phi$ 是态射，爆破后 $K+\Delta$ 的截面也会消失，故对的双有理等价不能只靠底簇的双有理性来定义。]

在使用 MMP 的大多数场合，需要在固定的簇上运行 MMP。它由以下步骤组成。

\par\medskip\noindent\textbf{3.31}（$\Q$-分解的 dlt 对情形的相对 MMP）\quad 从 $X$ 为域上的正规概形或正规复解析流形的对 $(X,\Delta)=(X_0,\Delta_0)$ 出发。设 $f_0:X\to S$ 是射影态射。

目标是确定造出中间的对 $(X_i,\Delta_i)$ 与射影态射 $f_i:X_i\to S$ 的递推过程。这个过程在重复若干次之后，应当以最终的对 $(X^*,\Delta^*)$ 与态射 $f^*:X^*\to S$ 结束。

\par\medskip\noindent\textbf{步骤 0（初始数据）}\quad 假定已经构成满足下列性质的对 $(X_i,\Delta_i)$ 与射影态射 $f_i:X_i\to S$。
\begin{enumerate}[label=（\arabic*）,leftmargin=2.8em,itemsep=1pt,topsep=2pt,parsep=0pt]
\item $X_i$ 是 $\Q$-分解的。
% [原文 p.124]
\item $(X_i,\Delta_i)$ 是 dlt 对。
\item $f_i$ 是射影的。
\end{enumerate}

\par\medskip\noindent\textbf{步骤 1（预备）}\quad 若 $K_{X_i}+\Delta_i$ 是 $f_i$-nef 的，则直接进入步骤 3(2)。若 $K_{X_i}+\Delta_i$ 不是 $f_i$-nef 的，则确立下面两个结果。
\begin{enumerate}[label=（\arabic*）,leftmargin=2.8em,itemsep=1pt,topsep=2pt,parsep=0pt]
\item （锥定理）$\NE(X_i/S)=\NE(X_i/S)_{K_{X_i}+\Delta_i\ge 0}+\sum\R^+[C_i]$。
\item （收缩定理）任意 $(K_{X_i}+\Delta_i)$-负的极射线 $R_i\subset\NE(X_i/S)$ 都可以收缩。设 $\mathrm{cont}_{R_i}:X_i\to Y_i$ 为对应的收缩（命题 2.5，定理 3.25），它给出下面的交换图式。
\end{enumerate}
\begin{center}
\begin{tikzcd}[column sep=4.5em, row sep=2.6em]
X_i \arrow[rr, "\mathrm{cont}_{R_i}"] \arrow[dr, "f_i"'] & & Y_i \arrow[dl, "g_i"] \\
{} & S & {}
\end{tikzcd}
\end{center}
% [Tutor Note: 步骤 1 把锥定理与收缩定理当作纲领的引擎：先分解曲线锥，再收缩负极射线；这一步依赖定理 3.25，而收缩后仍须保证 $(X_{i+1},\Delta_{i+1})$ 留在 dlt 对之类中。]

\par\medskip\noindent\textbf{步骤 2（双有理变换）}\quad 若 $\mathrm{cont}_{R_i}:X_i\to Y_i$ 是双有理的，则如下作出新的对 $(X_{i+1},\Delta_{i+1})$。
\begin{enumerate}[label=（\arabic*）,leftmargin=2.8em,itemsep=1pt,topsep=2pt,parsep=0pt]
\item （除子收缩）若 $\mathrm{cont}_{R_i}$ 是命题 2.5(2) 中的除子收缩，则令 $X_{i+1}:=Y_i$，$f_{i+1}:=g_i$，并令 $\Delta_{i+1}:=(\mathrm{cont}_{R_i})_*(\Delta_i)$。
\item （翻转收缩）若 $\mathrm{cont}_{R_i}$ 是定义 3.33 中的翻转收缩，则令 $(X_{i+1},\Delta_{i+1}):=(X_i^+,\Delta_i^+)$（$\mathrm{cont}_{R_i}$ 的翻转），$f_{i+1}:=g_i\circ f_i^+$。
\end{enumerate}
无论哪种情形，都要证明 $X_{i+1}$ 是 $\Q$-分解的、$f_{i+1}$ 是射影的、并且 $(X_{i+1},\Delta_{i+1})$ 是 dlt 对。此后，对 $(X_{i+1},\Delta_{i+1})$ 从步骤 0 重新开始。

\par\medskip\noindent\textbf{步骤 3（最终的结果）}\quad 结局是这个过程终止，并可望得到下列两种可能性之一。
\begin{enumerate}[label=（\arabic*）,leftmargin=2.8em,itemsep=1pt,topsep=2pt,parsep=0pt]
\item （Fano 收缩）若 $\mathrm{cont}_{R_i}$ 是命题 2.5(1) 中的纤维收缩，则令 $(X^*,\Delta^*):=(X_i,\Delta_i)$，$f^*:=f_i$。
\item （极小模型）若 $K_{X_i}+\Delta_i$ 是 $f_i$-nef 的，则同样令 $(X^*,\Delta^*):=(X_i,\Delta_i)$，$f^*:=f_i$。
\end{enumerate}
$(X^*,\Delta^*)$ 的特殊性质，大概会在试图通过 $(X^*,\Delta^*)$ 也理解 $(X,\Delta)$ 的尝试中得到使用。\hfill$\square$
